\nonstopmode{}
\documentclass[a4paper]{book}
\usepackage[times,inconsolata,hyper]{Rd}
\usepackage{makeidx}
\makeatletter\@ifl@t@r\fmtversion{2018/04/01}{}{\usepackage[utf8]{inputenc}}\makeatother
% \usepackage{graphicx} % @USE GRAPHICX@
\makeindex{}
\begin{document}
\chapter*{}
\begin{center}
{\textbf{\huge Package `regulationsdotgov'}}
\par\bigskip{\large \today}
\end{center}
\ifthenelse{\boolean{Rd@use@hyper}}{\hypersetup{pdftitle = {regulationsdotgov: Get Data from Regulations.gov}}}{}
\ifthenelse{\boolean{Rd@use@hyper}}{\hypersetup{pdfauthor = {Devin Judge-Lord; Maya Khuzam}}}{}
\begin{description}
\raggedright{}
\item[Title]\AsIs{Get Data from Regulations.gov}
\item[Version]\AsIs{0.0.0.9000}
\item[Description]\AsIs{This package provides functions to get data from regulations.gov.}
\item[License]\AsIs{MIT + file LICENSE}
\item[Encoding]\AsIs{UTF-8}
\item[Roxygen]\AsIs{list(markdown = TRUE)}
\item[RoxygenNote]\AsIs{7.3.2}
\item[URL]\AsIs{}\url{https://judgelord.github.io/regulationsdotgov/}\AsIs{}
\item[Imports]\AsIs{dplyr, here, httr, jsonlite, pdftools, purrr, quarto, stringr,
tibble}
\item[VignetteBuilder]\AsIs{quarto}
\item[NeedsCompilation]\AsIs{no}
\item[Author]\AsIs{Devin Judge-Lord [aut, cre] (<}\url{https://orcid.org/0000-0002-4384-4380}\AsIs{>),
Maya Khuzam [aut] (<}\url{https://orcid.org/0009-0007-7805-7460}\AsIs{>)}
\item[Maintainer]\AsIs{Devin Judge-Lord }\email{judgelor@umich.edu}\AsIs{}
\end{description}
\Rdcontents{Contents}
\HeaderA{get\_commentsOnId}{Get metadata for all comments on a document}{get.Rul.commentsOnId}
%
\begin{Description}
get\_commentsOnId("[objectId]") retrieves metadata for all comments on a document
\end{Description}
%
\begin{Usage}
\begin{verbatim}
get_commentsOnId(objectId, lastModifiedDate = Sys.time(), api_keys)
\end{verbatim}
\end{Usage}
%
\begin{Arguments}
\begin{ldescription}
\item[\code{objectId}] 
objectId obtained from a document's metadata 

\item[\code{lastModifiedDate}] 
Filter results by their last modified date. Default is set to \code{Sys.time()}. 

User-specified values must be formatted as \code{yyyy-MM-dd\%20HH:mm:ss}.


\item[\code{api\_keys}] API key(s) from api.data.gov. 

\end{ldescription}
\end{Arguments}
%
\begin{Examples}
\begin{ExampleCode}

# comment_metadata <- get_commentsOnId(commentOnId = "09000064865d514a", api_keys = "DEMO_KEY")

# head(comment_metadata)

# A tibble: 1 × 11
#   documentType      lastModifiedDate    highlightedContent withdrawn agencyId title objectId postedDate id   
#   <chr>             <chr>               <chr>              <lgl>     <chr>    <chr> <chr>    <chr>      <chr>
# 1 Public Submission 2024-07-26T13:29:3… ""                 FALSE     FBI      Comm… 0900006… 2024-07-2… FBI-…
# ℹ 2 more variables: lastpage <lgl>, commentOnId <chr>
\end{ExampleCode}
\end{Examples}
\HeaderA{get\_dockets}{Get dockets from an agency}{get.Rul.dockets}
%
\begin{Description}
get\_dockets("[agency\_acronym]") retrieves dockets for an agency
\end{Description}
%
\begin{Usage}
\begin{verbatim}
get_dockets(agency, lastModifiedDate, api_keys)
\end{verbatim}
\end{Usage}
%
\begin{Arguments}
\begin{ldescription}
\item[\code{agency}] 
Agency acronym (see official acronyms on regulations.gov)

\item[\code{lastModifiedDate}] 
Filter results by their last modified date. Default is set to \code{Sys.time()}. 

User-specified values must be formatted as \code{yyyy-MM-dd\%20HH:mm:ss}.


\item[\code{api\_keys}] API key(s) from api.data.gov. 

\end{ldescription}
\end{Arguments}
%
\begin{Examples}
\begin{ExampleCode}

# FBI_dockets <- get_dockets(agency = "FBI", api_keys = "DEMO_KEY") 

# head(FBI_dockets)

# A tibble: 6 × 8
#   docketType    lastModifiedDate     highlightedContent agencyId title                objectId id    lastpage
#   <chr>         <chr>                <lgl>              <chr>    <chr>                <chr>    <chr> <lgl>   
# 1 Rulemaking    2024-08-26T13:14:32Z NA                 FBI      "Bipartisan Safer C… 0b00006… FBI-… TRUE    
# 2 Rulemaking    2024-08-26T13:13:45Z NA                 FBI      "Child Protection I… 0b00006… FBI-… TRUE    
# 3 Rulemaking    2023-08-25T08:08:48Z NA                 FBI      "Recently Posted FB… 0b00006… FBI_… TRUE    
# 4 Rulemaking    2021-05-02T01:06:49Z NA                 FBI      "National Instant C… 0b00006… FBI-… TRUE    
# 5 Rulemaking    2021-05-02T01:06:49Z NA                 FBI      "FBI Criminal Justi… 0b00006… FBI-… TRUE    
# 6 Nonrulemaking 2021-05-02T01:06:49Z NA                 FBI      "FBI Records Manage… 0b00006… FBI-… TRUE 


##{DELETE?}
##---- Should be DIRECTLY executable !! ----
##-- ==>  Define data, use random,
##--	or standard data sets, see data().

## The function is currently defined as
#function (agency, lastModifiedDate = Sys.time()) 
#{
#    lastModifiedDate <- lastModifiedDate %>% ymd_hms() %>% with_tz(tzone = "America/New_York") %>% 
#        gsub(" ", "%20", .) %>% str_remove("\..*")
#    path <- make_path_dockets(agency, lastModifiedDate)
#    result <- purrr::map(path, GET)
#    status <<- tail(map(result, status_code), 1) %>% as.numeric()
#    url <- result[[20]][1]$url
#    if (status != 200) {
#        message(paste(format(Sys.time(), "%X"), "| Status", status, 
#            "| URL:", url))
#        Sys.sleep(6)
#    }
#    remaining <<- pluck(tail(map(result, headers), 1), 1, "x-ratelimit-remaining")
#    if (remaining < 2) {
#        message(paste(format(Sys.time(), "%X"), "- Hit rate limit, will continue after one minute"))
#        Sys.sleep(60)
#    }
#    docket_metadata <- purrr::map_if(result, ~status_code(.x) == 
#        200, ~fromJSON(rawToChar(.x$content)))
#    return(docket_metadata)
#  }
#
\end{ExampleCode}
\end{Examples}
\HeaderA{get\_documents}{Get documents from a docket folder}{get.Rul.documents}
%
\begin{Description}
get\_documents("[docketId]") retrieves metadata for all documents belonging to a docket 
\end{Description}
%
\begin{Usage}
\begin{verbatim}
get_documents(docketId, lastModifiedDate, api_keys)
\end{verbatim}
\end{Usage}
%
\begin{Arguments}
\begin{ldescription}
\item[\code{docketId}] 
ID of docket

\item[\code{lastModifiedDate}] 
Filter results by their last modified date. Default is set to \code{Sys.time()}. 

User-specified values must be formatted as \code{yyyy-MM-dd\%20HH:mm:ss}.


\item[\code{api\_keys}] API key(s) from api.data.gov. 

\end{ldescription}
\end{Arguments}
%
\begin{Examples}
\begin{ExampleCode}

# FBI_docket_2024 <- get_documents(docketId = "FBI-2024-0001", api_keys = "DEMO_KEY")

# head(FBI_docket_2024)

# A tibble: 1 × 16
#   documentType lastModifiedDate     highlightedContent frDocNum   withdrawn agencyId # commentEndDate     title
#   <chr>        <chr>                <chr>              <chr>      <lgl>     <chr>    <chr>       #        <chr>
# 1 Rule         2024-07-26T01:01:24Z ""                 2024-14253 FALSE     FBI      2024-08# -01T03:59:… Bipa…
# ℹ 8 more variables: postedDate <chr>, docketId <chr>, subtype <lgl>, commentStartDate <chr>,
#   openForComment <lgl>, objectId <chr>, id <chr>, lastpage <lgl>
\end{ExampleCode}
\end{Examples}
\HeaderA{get\_searchTerm}{Search for a term }{get.Rul.searchTerm}
%
\begin{Description}
get\_searchTerm("[searchTerm]") returns documents, comments, or dockets filtered by specific term
\end{Description}
%
\begin{Usage}
\begin{verbatim}
get_searchTerm(searchTerm, documents, lastModifiedDate, api_keys)
\end{verbatim}
\end{Usage}
%
\begin{Arguments}
\begin{ldescription}
\item[\code{searchTerm}] 
Term to filter by. 

\item[\code{documents}] 
Object to return. Default is set to \code{documents}. 

Valid options are \code{documents}, \code{comments}, or \code{dockets}. 


\item[\code{lastModifiedDate}] 
Filter results by their last modified date. Default is set to \code{Sys.time()}. 

User-specified values must be formatted as \code{yyyy-MM-dd\%20HH:mm:ss}.


\item[\code{api\_keys}] API key(s) from api.data.gov. 

\end{ldescription}
\end{Arguments}
%
\begin{Examples}
\begin{ExampleCode}
##---- Should be DIRECTLY executable !! ----
##-- ==>  Define data, use random,
##--	or do  help(data=index)  for the standard data sets.

## The function is currently defined as
function (x) 
{
  }
\end{ExampleCode}
\end{Examples}
\printindex{}
\end{document}
